\section{Ordered-pair descent and the genus-two saturation bound}
\label{sec:standard-ordered-pair-genus}

\paragraph{Status and provenance.}
HD-SS04--05, reserved question Q2ec, is separate from the frozen
HD-SS01--03.  The parent proposed the ordered-pair invariant argument
and the genus-two involution strategy.  The complete restriction,
field, splitting-degree and saturation arguments below were developed
by \texttt{/root/imprimitive}, with those contributions retained.
The complete author, root and independent whole-proof, source and
code readings pass without mathematical repair.  No claim that a character field automatically supplies a
model, or that the extra cohomology vanishes, is made.

We keep the actual finite coefficient setting of HD-SS01:
$V=W\otimes U$ is faithful and irreducible,
$W=\operatorname{Std}_n\circ x$, $A_n\subset x(J)\subset S_n$,
$n\ge4$, $b=\dim U$, and $H_1,H_2$ are the preimages of the
stabilizers of one letter and of an ordered pair of distinct letters.
Assume $V$ has its model over its exact imaginary quadratic field
$K$.  In particular $V\not\simeq V^*$.
Write $L=F_U=F_{U|H_1}$ as in HD-SS01.

\subsection*{HD-SS04. A lower-degree representation of exact field $K$}

\begin{proposition}[HD-SS04]
The representation $\tau_2=U|H_2$ is irreducible and nonselfdual,
and has exact character field
\[
                             F_{\tau_2}=K.
 \tag{SSX.1}
\]
It has a model over $L$.  If $s$ is its Schur index over $K$,
then $s\mid[L:K]$.  In particular an $L=K$ instance supplies a
$K$-model, but the general proposition does not assert $s=1$.
The actual induction identity is
\[
       \operatorname{Ind}_{H_2}^J\tau_2
                     =(W\otimes V)\oplus V.
 \tag{SSX.2}
\]
For a cover with $\mu_V=6$, the $\tau_2$-multiplicity variation
contains at least two copies of the $V$-variation and one copy of
the $U$-variation.  Thus it has complex dimension at least
$12+\mu_U$, rather than six.
\end{proposition}

\begin{proof}
Let $P$ be the actual permutation module on letters and $P_2$
the module on ordered pairs of distinct letters.  The tensor square
of $P$ has its diagonal and off-diagonal permutation summands,
and $P=1\oplus W$.  Consequently
\[
        P_2=P\otimes P-P=W\otimes W\oplus W.
 \tag{SSX.3}
\]
The displayed identity is one of actual semisimple representations;
it does not require separate irreducibility assertions for Young
diagram constituents, so it remains valid after restriction to
the actual alternating or symmetric quotient.
Since $A_n$ is transitive on ordered distinct pairs for $n\ge4$,
$P_2=\operatorname{Ind}_{H_2}^J1$.

The tensor projection formula yields (SSX.2).  For irreducibility,
using $W\simeq W^*$ gives
\[
 \begin{split}
 \dim\operatorname{End}_{H_2}(U|H_2)
  &=\dim\operatorname{Hom}_J(P_2,U\otimes U^*)\\
  &=\dim\operatorname{End}_J(W\otimes U)
       +\dim\operatorname{Hom}_J(U,W\otimes U)\\
  &=1+0=1.
 \end{split}
 \tag{SSX.4}
\]
The zero uses the different degrees $b$ and $(n-1)b$ of the
irreducibles $U$ and $V$.  Semisimplicity proves irreducibility
of $\tau_2$.

The same identity without a dual gives
\[
 \begin{split}
 \dim(U\otimes U)^{H_2}
   &=\dim(V\otimes V)^J
        +\dim\operatorname{Hom}_J(V^*,U)\\
   &=0+0=0.
 \end{split}
 \tag{SSX.5}
\]
The first zero is precisely the nonselfduality of $V$; the second
again follows from the degrees.  Therefore $\tau_2$ is nonselfdual.
For $g\in H_2$, its permutation fixes at least two letters.  If
$f(g)$ is the number of fixed letters, then
\[
              \chi_{\tau_2}(g)=\chi_U(g)
                    =\frac{\chi_V(g)}{f(g)-1}\in K.
 \tag{SSX.6}
\]
Hence $F_{\tau_2}\subset K$.  A finite-group irreducible character
is selfdual exactly when all its values are real.  Nonselfduality
therefore shows this field is not real, and the quadratic degree
of $K$ proves (SSX.1).

The $L$-model of $U$ from HD-SS01 restricts to an $L$-model of
$\tau_2$.  Regard it as a $K$-representation by restriction of
scalars.  After complexification it consists of $[L:K]$ copies
of $\tau_2$ at the chosen embedding of $K$, because its exact
character field is $K$.  The definition of the Schur index,
equivalently the simple-block description over $K$, says that
its multiplicity in a $K$-representation is divisible by $s$.
This is the explicit simple-block formula in Serre, Section~12.2,
pp.~92--93, read with its full written derivation
\cite{OG:Serre1977}.  It proves $s\mid[L:K]$.  It does not remove
the possible index using only the character-field calculation.

We also make the cohomological limitation explicit.  Put
$\tau_1=U|H_1$.  HD-SS01 proves it irreducible and shows it
occurs once in $V|H_1$.  Since
$W|H_1=1\oplus\operatorname{Std}_{n-1}$, this is a decomposition
\[
 V|H_1=\tau_1\oplus T,
 \qquad T=\operatorname{Std}_{n-1}\otimes\tau_1\ne0,
 \qquad\operatorname{Hom}_{H_1}(\tau_1,T)=0.
 \tag{SSX.7}
\]
Thus $\dim\operatorname{End}_{H_1}(V|H_1)\ge2$.
Since $\operatorname{Ind}_{H_1}^J1=1\oplus W$, reciprocity gives
\[
  \dim\operatorname{Hom}_J(V,W\otimes V)
   =\dim\operatorname{End}_{H_1}(V|H_1)-1\ge1.
 \tag{SSX.8}
\]
Also $U$ occurs exactly once in $W\otimes V$, because
$\dim\operatorname{Hom}_J(U,W\otimes V)
=\dim\operatorname{End}_J(V)=1$.
Equation (SSX.2) therefore contains at least two $V$ constituents
and one $U$ constituent, with every remaining multiplicity
nonnegative.  Apply the fixed rational or splitting-field
Frobenius correspondence to $H^1(C)$ as in HD-SS02.
It preserves all the summands and proves the claimed lower bound.
Passing to the faithful image of $\tau_2$ simply takes the
corresponding finite quotient of $C$; it does not discard any
of these cohomological constituents.
\end{proof}

\subsection*{HD-SS05. A quantitative genus-two exclusion}

\begin{theorem}[HD-SS05]
In the setting of HD-SS02, assume $b\ge2$, the original tuple
has at least twelve genuine entries, and the $V$-period rank is
nine.  If
\[
                              n>12b+36,
 \tag{SSX.9}
\]
then the natural letter quotient $X=C/H_1$ has genus at most one.
If it has genus one, its retained summand satisfies
$3\le\mu_U\le2b$.
The theorem does not assert $\mu_U=0$, and leaves the genus-zero
and genus-one cases open.
\end{theorem}

\begin{proof}
Use the nontrivial-$U$ identities from HD-SS01--02; exact-$K$
already ensures $U$ is nontrivial.  From (SS.17) and (SSX.9),
the nonnegative integer $\mu_U$ satisfies
\[
              \mu_U\le2b,
 \qquad
              \mu_U\ge(2b-3)(h-1)+3\quad(h\ge1).
 \tag{SSX.10}
\]
The first follows because $(4b+12)/(n-2)<1$.
For $h\ge3$ the second is at least $4b-3>2b$ for $b\ge2$,
which is impossible.  It remains to exclude $h=2$.

Suppose $h=2$.  The inequalities force $\mu_U=2b$.
Let $r_i=\operatorname{rank}(1-S_i)$, where $S_i=U(g_i)$.
The sphere Euler formula, the original sixfold formula,
and (SS.9) give
\[
 \sum_i r_i=4b,\qquad
 D=\sum_i\operatorname{rank}(1-V(g_i))=2(n-1)b+6,
 \qquad D'=6.
 \tag{SSX.11}
\]
In fact rank nine also forces exactly six genuine descended
marks, each of rank one, because a genus-two base with at most
five genuine marks has moduli dimension at most eight.
Only the equality $D'=6$ is needed in the counting below.

We now justify the involution reduction rather than assume it.
For any $x_i$, the multiplicity on $W(x_i)$ of a non-one root
$\lambda^{-1}$ of order $d\ge2$ is the number of permutation
cycles with length divisible by $d$, which is at most $n/d$.
The proof of (SS.17) assigned to every non-one eigenvalue of
$S_i$ a baseline moving contribution $(n-2)/2$.
All contributions above that baseline are nonnegative, including
the eigenvalue-one contribution.  If $S_i$ has an eigenvalue of
order $d\ge3$, even one such eigenvector contributes an additional
\[
 \left(n-1-\frac nd\right)-\frac{n-2}{2}
                       =\frac n2-\frac nd\ge\frac n6.
 \tag{SSX.12}
\]
The total available surplus over the baseline, by (SSX.11), is
\[
           D-\frac{n-2}{2}\sum_i r_i=2b+6.
 \tag{SSX.13}
\]
Condition (SSX.9) makes $n/6>2b+6$, so no such eigenvalue can
occur at any original inertia.  Every $S_i$ therefore has only
the eigenvalues $1,-1$ and, being of finite order, satisfies
$S_i^2=1$.

Let $t$ be the number of nonidentity $S_i$, and let $o_i$ be
the number of odd cycles of $x_i$.  The pullback operator
$S_i^{|c|}$ is $S_i$ on an odd cycle and the identity on an
even cycle.  Thus
\[
              \sum_{S_i\ne1}r_i o_i=D'=6,
 \qquad
              \sum_{S_i\ne1}o_i\le6.
 \tag{SSX.14}
\]
Every permutation with $o_i$ odd cycles has at most
$(n+o_i)/2$ cycles in total: each even cycle uses at least two
letters and each odd cycle at least one.  Its index is therefore
at least $(n-o_i)/2$.  Riemann--Hurwitz for the actual natural
quotient, which has genus two, now gives
\[
       2n+2=\sum_i\operatorname{ind}(x_i)
          \ge\frac{tn}{2}-3.
 \tag{SSX.15}
\]
As $n>12b+36>10$, this implies $t\le4$.
On the other hand $\sum_i r_i=4b$ and each nonidentity $S_i$
has $r_i\le b$.  Hence $t=4$ and all four ranks equal $b$.
Each such involution is $-I$.  The tuple generates $J$, so the
entire image of $U$ consists of scalars $\{1,-1\}$, contrary
to the irreducibility of $U$ in dimension $b\ge2$.

For completeness, the same counting prevents this exceptional
scalar configuration even if one had allowed $b=1$ without the
exact-$K$ assumption.  The four nontrivial $S_i$ then contribute
natural index at least $2n-3$ by (SSX.14).  At least eight other
genuine original entries remain, and their $S_i$ are the identity.
Genuineness forces each of their permutations to be nonidentity,
adding natural index at least eight.  This gives
$2n+2\ge2n+5$, a contradiction.  No assumption that twelve
descended or partner marks exist is used: these are the original
genuine coefficient entries.

Thus $h\le1$.  For $h=1$, (SSX.10) gives the asserted
$3\le\mu_U\le2b$.
\end{proof}
